% GDZ PPN236018647,native547middle–548before69;HKUST548–549.
\NumberedItem{64}
\indent Demonstrationes elegantiores pro nexu divisorum formae \(\square-a,+1\) cum \(-1,\pm2\) inveni.
\par\noindent\hfill[1797] Iun. 17.\quad Gotting[ae]

\medskip\begingroup\small
\indent Dies bezieht sich auf die in den artt. 147.—150. der \textit{Disquisitiones arithmeticae}, Werke I, S. 113 behandelte Frage. Es werden dort bezüglich der Teiler von \(x^2-\alpha\) die Fälle unterschieden:
\par\indent 1) \(\alpha\) von der Form \(4n+1\) oder \(-(4n-1)\),
\par\indent 2) \(\alpha\) von der Form \(-(4n+1)\) oder \(4n-1\),
\par\indent 3) \(\alpha\) von der Form \(\pm(4n+2)\).
\par\noindent Man würde also die Aufzeichnung zu lesen haben: Demonstrationes elegantiores pro nexu divisorum formae \(\square-\alpha\) cum \(+1,-1,\pm2\) inveni. (Nach einer Mitteilung von S. \textsc{Gundelfinger} in einem Briefe an F. \textsc{Klein} vom 10. April 1903.)
\par\noindent\hfill\textsc{Bachmann.}
\par\endgroup

\NumberedItem{65}
\indent Deductionem secundam theoriae polygonorum excolui.
\par\noindent\hfill[1797] Iul. 17.\quad Gotting[ae]

\medskip\begingroup\small
\indent Vergl. die folgende Notiz.
\par\endgroup

\NumberedItem{66}
\indent Per utranque methodum monstrari potest \textls[120]{puras} tantum aequationes solvi oportere.
\par\noindent\hfill[1797 Iul.]

\medskip\begingroup\small
\indent Außer der Methode des art. 352. der \textit{Disquisitiones arithmeticae} (Werke I, S. 431) hat \textsc{Gauss} für die Lösung der Kreisteilungsgleichung noch eine auf der Anwendung der \textsc{Lagrange}schen Resolvente (vergl.
\par\endgroup

\clearpage\SourceLeaf{548}\ArithmeticHead{519}{R}{1797 MART. 29—1797 IUL. 23.}
\begingroup\small
\noindent die Tagebuchnotiz Nr. 37) beruhende Methode angegeben, die im art. 360. der \textit{Disquisitiones arithmeticae} (Werke I, S. 450) auseinandergesetzt ist. Am Schluß des art. 18. der nachgelassenen Abhandlung \textit{Disquisitionum circa aequationes puras ulterior evolutio} (Werke II, S. 263) spricht \textsc{Gauss} von gewissen allgemeinen Untersuchungen, »quae theoriam \textls[120]{secundam} aequationum purarum in art. 360. \textit{Disquiss. Arithm.} inchoatam magis illustrant«; daraus geht hervor, daß er in der Tagebuchnotiz Nr. 65 mit der \textls[120]{deductio secunda} eben die Methode des art. 360. meint.\linebreak Es scheint übrigens, daß \textsc{Gauss} das Wesen dieser Methode schon erkannt hatte, als er die Tagebuchnotiz Nr. 55 schrieb, und daß es sich bei Nr. 65 nur mehr um die vollständige Durchführung handelte. Die Notiz Nr. 66 hebt hervor, daß beide Methoden die Lösung der Kreisteilungsgleichung durch reine Gleichungen, d. h. durch Wurzelziehen ermöglichen. Vergl. den Artikel 15 des Aufsatzes »Über \textsc{Gauss}' zahlentheoretische Arbeiten«, Werke X 2, S. 37.
\par\noindent\hfill\textsc{Bachmann.}
\par\endgroup

\NumberedItem{67}
\indent Quod Oct. 1. per ind[uctionem] invenimus demonstratione munivimus.
\par\noindent\hfill[1797] Iul. 20.

\medskip\begingroup\small
\indent Siehe oben S. 505 die Notiz Nr. 39 vom 1. Oktober 1796 und die zugehörige Anmerkung. Der Beweis findet sich im art. 358. der \textit{Disquisitiones arithmeticae}, Werke I, S. 445.
\par\noindent\hfill\textsc{Klein.\quad Bachmann.}
\par\endgroup

\NumberedItem{68}
\indent Casum singularem solutionis congruentiae
\[
x^n-1\equiv0
\]
\noindent (scilicet quando cong[ruentia] aux[iliaris] radices aequales habet), qui tam diu nos vexavit, felicissimo successu vicimus, ex congruentiarum solutione, si modulus est numeri primi \textls[120]{potestas}.
\par\noindent\hfill[1797] Iul. 21.

\medskip\begingroup\small
\indent Vergl. hierzu die \textit{Annotatio} zum art. 251. der \textit{Analysis Residuorum}, Werke II. S. 209, und die zu diesem Artikel gehörige Bemerkung von \textsc{Dedekind} ebenda, S. 241, ferner das Ende des art. 363. ebenda S. 232 und den art. 372. ebenda S. 237. Siehe auch den Artikel 16 des Aufsatzes »Über \textsc{Gauss}' zahlentheoretische Arbeiten«, Werke X 2, S. 42.
\par\noindent\hfill\textsc{Klein.\quad Bachmann.}
\par\endgroup
